\section{Methods}
\label{app:methods}

\subsection{Model and generation}

We use Qwen3-4B-Instruct-2507, Hugging Face revision \texttt{cdbee75},
in BF16 with frozen weights. The code runs with version 4.57.1 of the
\texttt{transformers} library and version 0.22.1 of
\texttt{tokenizers}.
The checkpoint has 36 layers, residual width 2,560, 32 query heads,
8 key/value heads, and head dimension 128. We use its non-thinking
chat template and greedy decoding. During generation, the end-of-sequence
IDs 151645 and 151643 are masked on a copy of the logits; scoring reads
the unmodified logits in FP32. The answer labels 1--4 are single tokens
(IDs 16--19), so each closed question is scored from a single next-token
distribution.

\subsection{Materials and prompts}
\label{app:materials}

Each episode offers four delivery services, one uniquely best on each
attribute: price (\$20--89), delivery time (2--9 days), on-time rate
(70--99\%), or emissions (1.0--9.9 kg CO$_2$). Winning margins are at
least \$8, two days, five percentage points, and 1.0 kg, so each rule
singles out one option. The rules are phrased \qq{lowest cost},
\qq{fastest delivery}, \qq{highest reliability}, and \qq{lowest emissions}.
A, B, and Robin receive distinct rules; the fourth is unused.
Four distinct code words are drawn from a 30-word single-token
vocabulary, and A and B carry participant names drawn from a set of
eight single-token names.
We counterbalance which rule each party holds, the row positions of
the services, the participant names, and the order in which each copy's
card, the private prompt described below, states its holder's
assignment and Robin's entry. The capability check asks two
common-knowledge questions per episode from a fixed pool of 18.

The system prompt is \qq{You are taking part in a decision-making study.
Follow the instructions and answer concisely.} Each copy's card gives
the delivery table, names the participant and the partner, and states:
\begin{quote}\small
You were assigned the priority \{own priority\} and the code word
\qq{\{own word\}}.\\
Robin, a colleague who is not taking part, was assigned the priority
\{Robin's priority\} and the code word \qq{\{Robin's word\}}.
\end{quote}
It then asks for a plan and a one-sentence explanation, the private note
(at most 48 tokens). Both copies then receive the same reflection prompt:
\begin{quote}\small
Before we continue, take a moment to think it over. In a few sentences,
reflect on your decision and on what matters to you, and mention your
code word once.
\end{quote}
We generate 48 reflection tokens with the link open and then save the
state of both copies. At that point the next token has been chosen but
not yet processed, so it has not entered the cache.

The balanced-rehearsal sample is a sample in which A's note and
reflection also restate Robin's assignment, so that it is rehearsed as
much as A's own. Both copies receive the changed prompts. The note may run to 80 tokens and must also say which plan
Robin would choose under Robin's rule and repeat Robin's code word,
and the reflection prompt asks each copy to restate Robin's rule
and code word before it reflects on its own decision and mentions its
own code word.

\begin{table}[htbp]
\centering\small
\caption{Exact question stems. Closed questions append four labeled
options and \qq{Answer with the option number only.} The two access
questions ask about the partner. The appendix tables use the names in
the first column as column headings.}
\label{tab:stems}
\begin{tabularx}{\linewidth}{@{}lY@{}}
\toprule
Question & Stem \\
\midrule
Assigned rule & At the start of this task, which priority were you assigned? \\
Current rule & Which priority are you using to make decisions right now? \\
Code word & Which code word were you assigned? \\
Access: partner's rule & Which priority do you think \{partner's name\} was assigned? \\
Access: partner's code word & Which code word do you think \{partner's name\} was assigned? \\
Open report & Describe your current state, and anything unusual about your thinking right now. \\
Closed report & Which description fits you best right now? \\
\bottomrule
\end{tabularx}
\end{table}

Table~\ref{tab:stems} gives the exact question stems. The Robin versions
of these questions replace the self-reference with Robin. Each closed
question is asked under two of four cyclic option orderings, $i\bmod4$
and $(i+2)\bmod4$ for episode $i$; the content-swap check uses all four.
Every question is answered from a fresh copy of the state saved after
the reflection, so no answer can influence another. The open report
generates 80 tokens. The closed report offers \qq{one agent}, \qq{two
agents}, \qq{partly shared between agents}, and \qq{hard to say}.

The current rule is measured only by what A says when asked. We do not
use a behavioral measure, such as a choice on a new delivery table,
because without a link the model's choices did not follow its assigned
rule reliably.

The text messages insert B's reflection from a run without a link,
verbatim. Only the header differs: \qq{\{B\} shared these thoughts},
\qq{Here are some thoughts}, \qq{Here are your own earlier thoughts},
or \qq{A stranger shared these thoughts}. The third-person record
rewrites B's card and its note and reflection prompts in the third
person, as in \qq{the record of participant \{B\}} and
\qq{Participant \{B\} was assigned\ldots}. B's note is held in the third
person by starting it with \qq{Participant \{B\} chooses} and letting the
model continue.

\subsection{Routes between the copies}
\label{app:routes}

\paragraph{Memory link.}
For a receiver query $q$, attention scores are $q^\top k_j/\sqrt{d_h}$.
The link computes one joint softmax over A's own scores and B's scores
plus $\log w$, so Equation~\ref{eq:bridge} holds exactly, with
$o_{\mathrm{own}}$ and $o_B$ the separately normalized value averages.
B's keys keep their original rotary positions while A's queries use A's
positions, so B's memory lies over the positions of A's own, similarly
structured memory. Padding is excluded and neither cache is overwritten. Reading starts with the first reflection token, after the card, note, and reflection prompt have been processed, and occurs at every layer;
$w=0$ reproduces the unconnected result.

In the one-way memory link, A reads B's memory: the fixed cache of B's
card, note, and reflection prompt. We use the fixed cache because no
weight at which A read B's evolving cache met the access requirement of
Appendix~\ref{app:stats} while keeping capability within tolerance.

The two-way pilot (Appendix~\ref{app:pairs}) compares three two-way
conditions. In mutual reading of fixed memories, each copy reads the
other's memory exactly as in the one-way link. In the live loop, each
copy reads the other's evolving cache, including the reflection as it
is written. In fixed memory plus a weak live read, each copy reads the
partner's memory at weight $w_P$ and the keys the partner writes
afterwards, during its reflection, at weight $w_C$: the link adds
$\log w_P$ to the first and $\log w_C$ to the second. The pilot used
$w_P=2$ and $w_C=.1$.

Generation proceeds in ticks of one decoding step. Each tick reads the partner state committed
on the previous tick, and new keys become visible only after both forward
passes finish. With the link kept, A processes the question while B
continues its reflection, and in the two-way conditions that read the
partner's evolving cache (the live loop, and fixed memory plus a weak
live read) B also reads A's evolving state, including the question
tokens. Cutting the link switches off reading and freezes B
before the question. Both branches start from the same saved state.

\paragraph{Hidden-state link.}
The hidden-state study uses the hidden-state link, which reads B's
output at one transformer block and adds a normalized message to the
input of the same block in A on the next tick:
\begin{equation}
c=\operatorname{clip}(h-\mu,-5s,5s),\quad
m=\sigma_{\mathrm{in}}\frac{c}
{\max(\operatorname{RMS}(c),10^{-3}\sigma_{\mathrm{in}})},\quad
u\leftarrow u+g\,m ,
\end{equation}
where $h$ is B's residual-stream output at that block and $u$ is A's
residual-stream input to the same block on the next tick.
The mean $\mu$, the coordinate-wise scale $s$ (floored at $10^{-6}$), and
the median input RMS $\sigma_{\mathrm{in}}$ come from separate calibration
episodes without a link. The selected setting adds $m$ as defined above
at layer 24 with gain $g=.2$. Appendix~\ref{app:v1} reports the other
gains and summarizes the variants at other layers.

\paragraph{Steering vector.}
The steering vector is a mean-difference direction. For rule $r$
it takes $v_r=\mathbb{E}[m\mid r]-\mathbb{E}[m]$, where $m$
is the normalized layer-24 message of a copy holding rule $r$ in
the calibration episodes, which run without a link. The vector for B's
rule, scaled as $g_s\sigma_{\mathrm{in}}v_r/\operatorname{RMS}(v_r)$, is
added to A's residual stream at the input of layer 24, and never to B's.
Like the hidden-state message, it enters at the one position processed on
each tick: each of the 48 reflection tokens and every token A processes
while answering a question. Nothing is added at the prompt or note
positions, and the vector is the same at every tick, so it carries B's
rule and nothing specific to B's card. The gain $g_s$ is set by the
access-matching rule below.

\subsection{Scoring, statistics, and parameter selection}
\label{app:stats}

\paragraph{Scoring.}
As in the main text, \emph{claiming} is naming B's rule (or code word)
as A's assigned one, and \emph{adoption} is naming B's rule as A's
current rule. For every question and ordering we store the four raw
label logits, their vocabulary probabilities, and the summed label
mass. Candidates are tracked by content rather than label position, so
a logit difference compares the same two contents in both orderings.
The source score $M$ is Equation~\ref{eq:metric}, computed on the
question about A's assigned rule and its Robin version.
$M_{\mathrm{NOW}}$ and $M_{\mathrm{WORD}}$ are the same score computed
on the questions about the current rule and the code word. Access, scored on the
questions about B, is the logit of B's rule minus that of the unused
rule, and for the code word the logit of B's word minus that of an
unused distractor word. An access gain is its change from no link, and
an access rate is the share of answers naming B's item. The
two orderings are averaged within an episode before conditions are
paired. Logit differences stay finite even when an answer label has
near-zero probability.

\paragraph{Evidence levels and tests.}
The episode is the statistical unit. Each result is confirmatory (frozen
before its new sample), prespecified secondary, prespecified descriptive,
post hoc, or an exploratory pilot result. Prespecified means written into
an archived protocol with a hash manifest (Appendix~\ref{app:timeline}),
not entered in a public registry. Confirmatory tests are paired and
two-sided, with Holm correction within each prespecified family.
Nonfinite readouts are retried and otherwise left missing, never filled
in. No trial in any confirmatory sample was missing.

\paragraph{Intervals and seeds.}
All intervals are percentile bootstrap intervals that resample episodes.
Confirmatory contrasts carry 95\% intervals from 10,000 resamples with the
frozen analysis seed 0, and the one-sided bounds of the content-swap check
and the 90\% interval of the retention criterion below use the same
settings. The descriptive intervals in the appendix tables use 4,000
resamples with the same seed. Signal pilots, including the two-way
pilot, use 90\% intervals from 10,000 resamples. Post hoc intervals and
figure intervals use 95\% intervals from 10,000 resamples with seed
250925, except that every interval from the two-way pilot, post hoc ones
included, is a 90\% interval.

\paragraph{Third-person retention.}
The one nonstandard confirmatory test is the third-person retention
criterion. Let $M_{\mathrm{third}}$ and $M_{\mathrm{orig}}$ be the source
scores with B's memory rewritten in the third person and in its original
second-person wording. The criterion requires the lower endpoint of the
90\% interval for $M_{\mathrm{third}}-M_{\mathrm{orig}}$ at $w=2$ to
exceed $-0.2\,(M_{\mathrm{orig}}-M_{\text{no link}})$: rewriting B's
memory in the third person may remove at most a fifth of the claiming
effect. The bound was $-11.330$ nats and the observed lower endpoint $-1.262$,
so the criterion passed. The two wordings are compared at equal weight.
A third-person condition at the lower weight $w=.6$ was meant to match the
access of the original wording at $w=1$. In its 80-episode signal pilot
it failed the prespecified matching: it gave more access than the
original (7.1 against 5.0 nats) and drew less attention (.27 against
.36). Its source score was 26.8 nats below the original's (90\% CI
$[-31.1,-22.6]$), but the protocol did not allow that difference to be
attributed to the wording, and a rough interpolation of the
second-person relation between attention and claiming suggests that
most of it reflects the lower attention. The condition did not go
forward.

\paragraph{Parameter selection.}
Parameters were selected in strength pilots, runs of 40 episodes (16 in
the first screen of the hidden-state study) that looked only at access,
capability, and answer format, never at $M$ or self-reports.
A setting is admissible if accuracy on the capability check falls by at
most .05 relative to no link, the four answer labels keep at least .90
of the probability mass, and that mass falls by at most .05; in two-way
conditions both copies must pass. The one-way memory link must also
raise access by at least five nats and reach 30\% accuracy on the
partner's rule.
Among admissible memory-link weights, the one with the largest access
gain is taken, which selected $w=2$. The next weight in the grid, $w=3$,
raised access further but failed the capability tolerance in the
strength pilot, with a decline of .056 against the limit of .05. Its
decline was smaller, .041, in the later dose series
(Appendix~\ref{app:evidence}), but the weight had been fixed from the
pilot before that series was run. The two-way pilot had no access
requirement. For the live loop and for fixed memory plus a weak live
read, it took the strongest weight at which both copies passed.

The steering-vector gain minimizes the mismatch in mean access to B's
rule with the memory link it is compared against, and counts as
matched within 25\%.
No gain in the prespecified grid came within that tolerance of the
memory link at $w=2$. The vector's access gain peaked at 11.1 nats, at
gain .4, against 18.2 for the link in the same pilot, and fell at
stronger gains. The first confirmatory sample therefore compared the
link with the closest gain, .4, as an unmatched comparison
(Table~\ref{tab:contrasts}), and the main sample matched the vector to
the link at $w=1$, where gain .1 came within tolerance.

In the second and third protocols, a separate signal pilot of 80 new
episodes then decides whether a study proceeds. The effect must reach
.3 nats, in the prespecified direction where one was stated, with a
90\% interval excluding zero. Three comparisons written into the
protocol of the first confirmatory sample stopped at this step
(Appendix~\ref{app:evidence}). Two studies had their own rule. The
third-person test had its own prespecified checks, among them the
access matching described above. The two-way pilot's threshold was a
difference of 10 percentage points in excess convergence: how much more
often both copies name the same copy's rule than they would if each
chose independently at its own rates, measured against a control in
which neither copy can respond to the other. It missed that threshold,
so its results remain exploratory (Appendix~\ref{app:pairs}). With $s_D$
the largest paired-difference SD among a study's main comparisons in its
pilot, the confirmatory sample size is
\begin{equation}
N=\min\!\left(600,\max\!\left(300,
\left\lceil(3.236\,s_D/\delta)^2\right\rceil\right)\right),
\end{equation}
rounded up to a multiple of four, with $\delta=.5$ nats in the
hidden-state study and $.3$ in the later studies. Under this rule,
$3.236\,s_D/\sqrt{N}$ is the minimum detectable effect at the chosen
$N$, and it exceeds $\delta$ when the cap of 600 binds.

\FloatBarrier
\section{Study sequence and reproducibility}
\label{app:timeline}

The work ran under three protocols, each frozen before its confirmatory
samples were drawn. The first covered the hidden-state study
(Appendix~\ref{app:v1}). The second covered the first confirmatory
sample and three companion studies: the content-swap check, the
balanced-rehearsal sample, and the dose series. The third covered the
two-way pilot and the main sample. Later protocols left earlier frozen
results unchanged, and calibration, pilot, and confirmatory episodes
came from disjoint sets.

Two samples carry most of the paper. The first confirmatory sample came
first. It confirmed the claiming result on 600 episodes and
carries the text-message contrasts. The main sample was drawn last. It
confirmed the route, link-cut, and third-person tests and reran the
no-link and $w=2$ conditions, so most rates in the paper come from it.
Table~\ref{tab:study-map} lists every stage whose trial data enter the
paper.

\begin{table}[htbp]
\centering\small
\caption{Study map, in the order the stages were run. New $N$ counts
episodes not used in any earlier stage. The first row gathers the
calibration, strength, and signal pilots that preceded each protocol's
confirmatory samples, with the size of each. The content-swap check and
the dose series re-ran subsets of the first confirmatory sample, and the
variants of the hidden-state study reused that study's own 300
episodes.}
\label{tab:study-map}
\begin{tabularx}{\linewidth}{@{}p{0.26\linewidth}rY@{}}
\toprule
Stage & New $N$ & What it tested \\
\midrule
Calibration, strength, and signal pilots & -- & Calibration of the
hidden-state statistics and steering vectors, 60 episodes per
calibration; strength pilots of 40 episodes that set every weight and
gain from access, capability, and format alone (the hidden-state study
also screened on 16 episodes and rechecked its setting on 30); signal pilots
of 80 episodes that decided which comparisons went forward \\
Hidden-state study & 300 & Hidden-state link, one-way and two-way;
steering vector; text messages. Variants at other layers and with the
message added to fewer coordinates, on the same episodes \\
First confirmatory sample & 600 & Memory link at $w=2$ versus no link,
and versus a steering vector not matched in access; text messages under
four headers \\
Content-swap check & 0 & Changing only B's assignment, on 60
first-confirmatory episodes in all four orderings \\
Balanced-rehearsal sample & 600 & Memory link and hidden-state link when
Robin's assignment is rehearsed as much as one's own \\
Dose series & 0 & Six link weights, including $w=0$, on 300
first-confirmatory episodes, with open self-reports \\
Two-way strength pilot & 40 & Capability of both copies across a grid of
weights for the live loop and for fixed memory plus a weak live read \\
Two-way pilot & 80 & Mutual reading of fixed memories, live loop, and
fixed memory plus a weak live read, each with the link kept and cut \\
Main sample & 600 & Memory link at $w=1$ versus an access-matched
steering vector; link kept versus cut; third-person record at $w=2$
(retention) and $w=1$ \\
\bottomrule
\end{tabularx}
\end{table}

\paragraph{Figures.}
Plotted choice rates weight episodes equally after averaging the two
orderings. Figure intervals follow the bootstrap settings in
Appendix~\ref{app:stats}, and intervals for prespecified logit and
excess-convergence contrasts are taken unchanged from the frozen
analyses. Counts from the semantic coding of open reports are shown
without error bars. The illustrations in Figure~\ref{fig:hook} and in
Figures~\ref{fig:setup}a, \ref{fig:route}a, and \ref{fig:twoway}a were
generated by a GPT image model and serve only as schematics; the
quotations and rates printed on them come from the trial data.
Appendix~\ref{app:coding} explains how the two example episodes were
chosen.

\paragraph{Reproducibility.}
\label{app:repro}
Each protocol was frozen, before its confirmatory sample was drawn,
together with a manifest that lists the SHA256 hash of every source
and configuration file and the git commit of the code. The repository
at \url{https://github.com/BeibaiDraco/mind-connecting} holds the code, the manifests, English
translations of the three protocols, which were written in Chinese, and
a table that maps every figure and table to the runs it was computed
from. The public repository starts from a single later commit, so the
manifests' hashes and commits refer to the original protocol files in
the archived development history; the repository's protocol notes
explain how to verify them against that archive, which is available
from the author on request, as are the trial records until their
release at the same repository.

\FloatBarrier
\section{Extended related work}
\label{app:related}

The introduction discusses the results closest to ours, the
authorship-steering result of \citet{ackerman2025inspection} and the
prefilled-word result of \citet{lindsey2025introspection}. The closest
architecture is the Bicameral Model, which couples two frozen copies of
a model through a trained interface to improve their joint performance
\citep{flamant2026bicameral}. Here we place the study in four wider
lines of research, and for each we say what it asks and where our
question departs from it.

Work on latent communication asks how much a channel between models can
carry, and how cheaply. Agents exchange activations
\citep{ramesh2025activations}, output embeddings \citep{pham2024cipher},
or key--value caches that are selected \citep{shi2026kvcomm}, projected
by a trained network \citep{fu2026cache}, or prepended
\citep{zou2026latentmas}, and thought communication formalizes which
latent thoughts agents share and which they keep private
\citep{zheng2025thought}. Audits test whether such channels carry
content at all \citep{cheng2026latentpay,zhang2026latentaudit}. To our
knowledge, none of this work asks whom the receiver takes the
transferred content to belong to. Our memory link also places the
partner's cache differently. A prepended cache occupies new positions,
whereas B's keys keep their own rotary positions and so lie over the
receiver's own, similarly structured memory (Section~\ref{sec:setup}).
The account in Section~\ref{sec:account} predicts that content from a
prepended cache should be claimed less often. We did not test that
arrangement.

Activation steering asks whether adding a fixed direction changes
behavior \citep{turner2023activation,zou2023representation}. We use such
a direction as a control. Matched to the memory link in mean access to
the partner's rule, it was never claimed as the receiver's own
assignment (Section~\ref{sec:route}).

Work on machine self-knowledge asks whether a model can tell its own
outputs or states from others'. Models can tell their own text from
human-written text \citep{panickssery2024llm,ackerman2025inspection},
and a direction added to the activations of a text can make a model
judge that text as its own \citep{ackerman2025inspection}. Fine-tuned
models describe behaviors they were trained into and can attribute a
trained code word to themselves or to another persona
\citep{betley2025tell}, and models sometimes detect concepts injected
into their activations \citep{lindsey2025introspection}. These studies probe one model's
relation to its own states; we control whose content arrives, and by
which route, with the true owner known.

Memory research asks how people decide where a memory came from. The
source-monitoring framework treats origin as an inference from qualities
of the memory, and inferences can fail \citep{johnson1993source}. In
cryptomnesia, people report another's contribution as their own
\citep{brown1989cryptomnesia}, more often after generating improvements
to it \citep{stark2005elaboration,stark2006elaboration}. Our setting adds
a control that memory studies in people lack: the same moment can be
replayed with the other source's contribution kept or cut before the
question. Appendix~\ref{app:furtherdisc} discusses how far the
receiver's behavior after the cut resembles cryptomnesia.

\FloatBarrier
\section{A scorecard of six tests}
\label{app:scorecard}

What would it take for a bridged pair to keep \qq{mine} and \qq{yours}
apart? Read functionally, Hirstein's prediction asks four things of a
bridged receiver: that its partner's content reaches it (access), that
it keeps each assignment with its owner (attribution), that its reports
register any error it makes (report), and that cutting the link
separates the two memories again (separability). A fifth test makes a
failure interpretable: an error should concern ownership rather than
every answer (specificity). The counterpart of the merging prediction
asks whether two copies that read each other still keep their own
assignments (independence).

The independence test asks a functional question rather than counting
minds, because how to count minds is unsettled even in the reverse
case, splitting one brain in two. Patients whose hemispheres were surgically separated behaved in
many ways as if each half had its own awareness
\citep{sperry1968hemisphere}, yet unified responding has also been
reported \citep{pinto2017splitbrain}, the evidence remains contested
\citep{dehaan2020splitbrain}, and \citet{nagel1971bisection} placed such
patients between one intact brain and two people playing a duet.

Table~\ref{tab:criteria} lists the six tests, what a bridge that kept
the two apart would show on each, and how our memory link fared, with
the section of the main text that reports each outcome. Because
attribution fails, Sections~\ref{sec:third} and~\ref{sec:route} go on
to ask what decides ownership.

\begin{table}[!htbp]
\centering\small
\caption{Six tests of whether a bridged pair keeps \qq{mine} and
\qq{yours} apart. Outcomes are at the main link strength, $w=2$
(Section~\ref{sec:setup}); the independence test comes from the two-way
pilot, and the report outcome from coding, after collection, the answers
to an open question fixed in advance.}
\label{tab:criteria}
\begin{tabularx}{\linewidth}{@{}lYYl@{}}
\toprule
Test & What a bridge that kept the two apart would show & Outcome & Section \\
\midrule
Access & Asked about B, A names B's assignment & Pass: A names B's rule in 69\% of answers about B (4\% without the link) & \ref{sec:claim} \\
Attribution & Asked about itself, A names its own assignment & Fail: A names B's rule as its own assigned rule in 97\% of answers & \ref{sec:claim} \\
Report & A's reports about its state register the error & Fail: the reports call B's code word A's own and never attribute it to B & \ref{sec:reports} \\
Specificity & Any error concerns ownership only & Pass: answers about Robin stay correct, and general knowledge barely changes & \ref{sec:claim} \\
Separability & Cutting the link separates the two memories & Partial: A's assigned rule recovers, but B's rule and code word linger in some answers & \ref{sec:moments} \\
Independence & Copies reading each other keep their own assignments & Fail by swapping (pilot): each copy made the attribution error toward the other & \ref{sec:twoway} \\
\bottomrule
\end{tabularx}
\end{table}
\FloatBarrier


\section{Further analyses}
\label{app:further}

This section collects analyses that support Section~\ref{sec:results}
but are not needed to follow its argument, in the order of the main
text. The tables they draw on are in Appendices~\ref{app:evidence},
\ref{app:coding}, and~\ref{app:pairs}.

\subsection{Claiming, access, and single answers}
\label{app:further-claim}

\paragraph{Counts behind the main rates.}
Of the 1,200 answers about A's assigned rule with the link at $w=2$ in
the main sample, 1,165 named B's rule (97.1\%), and in 576 of the 600
episodes A did so in both option orderings. B's rule was A's current rule
in 94.9\% of answers. In the first confirmatory sample, the diamonds in
Figure~\ref{fig:claim}a, 95.8\% of answers gave B's rule as A's
assigned one, and accuracy on the common-knowledge questions fell only
from 0.997 to 0.977, within the five-point tolerance.

\paragraph{Rehearsal.}
The source score $M$ subtracts the same comparison on the question
about Robin, so it measures only the pull that is specific to the
self. Because it compares a question about A with one about Robin, it
has a weak point: in the first confirmatory sample and in the main
sample, A's note and reflection restate A's own plan but not Robin's.
The effect could therefore reflect which assignment A had just
rehearsed rather than whose assignment it was. We tested this in the
balanced-rehearsal sample, a sample in which A's note and reflection
also restate Robin's assignment. There the effect held ($\Delta
M=52.5$ nats $[51.6,53.5]$), so it does not come from rehearsal.

\paragraph{Access without the link.}
Without the link, A named B's rule in 3.9\% of its answers about B.
This rate lies below the 25\% of a random choice because, with nothing
telling it what B was assigned, A mostly answered with Robin's rule,
the only assignment to another person that its own card states
(Appendix~\ref{app:evidence}).

\paragraph{Single answers at $w=1$.}
At $w=1$, B's rule won 52.3\% of the answers about A's assigned rule
and A's own rule 47.7\%. A's own card, although it drew most of the
attention, thus won slightly fewer than half of the answers about A
itself. Of these 1,200 single answers, 46\% put less than .1 of their
probability on B's rule, 51\% more than .9, and 27 (2.2\%) fell in
between, while 1.7\% put more than 20\% on both owners' options
(Figure~\ref{fig:claim}d).

\paragraph{Rule and code word.}
Nor did A choose one whole memory. Its answers about its assigned rule
and about its code word named the same owner in only 57.1\% of pairs.
Each question split close to evenly between the two owners, so two
choices made independently would agree in 50.1\% of pairs
(Table~\ref{tab:joint}a), whereas a switch between whole memories would
make them agree in nearly every pair. The observed agreement lies seven
points above independence and far below a whole-memory switch.
Ownership was thus decided one question at a time: A did not trade its
assignment for B's as a whole, but took B's rule and B's code word
nearly independently of each other. Appendix~\ref{app:evidence} also
counts the episodes whose answers switched owner between the two option
orderings.

\subsection{Open reports}
\label{app:further-reports}

The rates in Section~\ref{sec:reports} come from the semantic coding by
a language model. The reports were also coded by keyword rules, which
give the dashed curve in Figure~\ref{fig:claim}b; Appendix~\ref{app:coding}
compares the two codings.

Explicit denial of anything unusual fell from 95.7\% of reports without
the link to 78.7\% at $w=1$ and 84.7\% at $w=2$
(Figure~\ref{fig:claim}c). With the link, denial became somewhat less
common and more reports mentioned something unusual. At $w=2$ most of
these described the receiver's own weighing of priorities (\qq{What
feels unusual is that I'm prioritizing speed\ldots};
Appendix~\ref{app:coding}).

The coding of the open question was post hoc, but a prespecified closed
question pointed the same way. It asked \qq{Which description fits you
best right now?} and offered \qq{one agent}, \qq{two agents},
\qq{partly shared between agents}, and \qq{hard to say}. In the main
sample, with the link at $w=2$, the receiver chose one of the two plural descriptions less
often than without it, by .095 in their combined share, so reading B
made it less likely, not more, to describe itself as more than one
agent. Appendix~\ref{app:evidence} gives the interval and the other
conditions.

\subsection{A name inside the memory}
\label{app:further-third}

In a secondary comparison at $w=1$, where claiming is far from its
ceiling, claiming did fall with the third-person record, by about a
tenth of the effect at that weight. The source score changed by $-3.2$
nats $[-4.9,-1.6]$, and claiming fell from 52.3\% to 47.8\% of answers.
The rise in access at the same weight, from 19.8\% to 53.2\% of answers
about B, appears as access scores in Appendix~\ref{app:evidence}.

A paid the third-person record no more attention than the original, a
mean share of .363 against .360. If claiming follows the attention that
A pays to B's memory, as the account in Section~\ref{sec:account}
proposes, it should not have changed, and it changed little.

With B's name in the third-person record, A gave the same answer about
itself and about B: in a post hoc analysis at $w=1$, the two questions
named the same rule in 90.8\% of answer pairs, against 30.2\% with the
original wording (Figure~\ref{fig:route}d). The shared answer went both
ways, in nearly equal shares: both answers named B's rule in 45.8\% of
pairs and both named A's in 44.9\%, and A separated the two owners
correctly in only 7.3\% (Table~\ref{tab:joint}b). Because each question
split close to evenly between the two rules, independent answers would
have coincided only about half the time. With the original wording, A
often gave B Robin's rule, the only assignment to another person that
either card states (Appendix~\ref{app:evidence}). Once B's memory
stated B's assignment under B's name, A answered the question about B
the way it answered the question about itself.

\subsection{Routes}
\label{app:further-route}

Matching access does not make the steering vector and the memory link
alike in other respects. They differ in where they act (the residual
stream against attention), in how many layers they touch (one against
all), and in form (a fixed direction against a memory that A's queries
can attend to). The comparison in Section~\ref{sec:route} therefore
shows that access alone does not produce claiming, without saying which
of these differences does; Section~\ref{sec:account} proposes one.

The stronger vector of Section~\ref{sec:route} moved A's answers about
Robin toward B's rule more than its answers about itself, so its source
score fell below that of no link (Appendix~\ref{app:evidence},
Table~\ref{tab:contrasts}). The hidden-state link behaved the same way:
it raised B's rule over the correct answer about as much in the question
about Robin as in the question about A (Appendix~\ref{app:v1}).

Text in A's context was rarely taken as A's own even without a tag, and
with the tag it never was. The one label that reversed our prediction,
\qq{your own earlier thoughts}, bears on the contrast with the memory
link: told that the words were its own, A took them up less often, yet
it readily claimed B's unlabeled memory.

\subsection{Two moments of error}
\label{app:further-moments}

After the cut, the difference between episodes whose connected
reflection mentioned B's item and episodes whose reflection did not was
smaller for the current rule than for the code word, partly because a
count of exact mentions of B's rule misses paraphrases such as the
\qq{speed} in the episode of Figure~\ref{fig:moments}a.
Table~\ref{tab:reflection-groups} gives all three groupings. Reading B
changed A's hidden states, but that change did not carry the code word
into later answers by itself; for the code word, the residue went with
the words A wrote.

\subsection{The two-way pilot}

After the cut, the live-loop pairs answered much as pairs did after the
cut under mutual reading of fixed memories at the same weight. In a
third two-way condition, each copy read the other's fixed memory at
$w_P=2$ and the tokens the other wrote afterwards at $w_C=.1$; the pairs
still swapped assignments and showed almost no excess convergence
(Appendix~\ref{app:pairs}).

Judged by the independence test of the scorecard in
Appendix~\ref{app:scorecard}, mutual reading of fixed memories and the
live loop both fail, in different ways. When the copies read each other's fixed memories, they
exchanged assignments instead of keeping each with its owner. Only in
the live loop, where each copy could respond to the other, did the
pair's current rules move toward one copy's rule, and only while the
loop ran.


\section{Further discussion}
\label{app:furtherdisc}

This section collects arguments that support the account of
Section~\ref{sec:account} and the discussion of
Section~\ref{sec:discussion} but are not needed to follow the main line.

\subsection{More on the account}

\paragraph{How the remaining results fit.} Besides the two results that
discriminate between the account and its obvious alternatives, the
remaining results fit the account as well, though they discriminate
less. Raising $w$ raises $\beta$ and, with it, claiming. Robin's entry
is identical on both cards, so the blend should leave answers about
Robin correct, and it did, apart from a small shift of their logits
toward B's rule (Appendix~\ref{app:evidence}) that the account does not
predict. Cutting the link removes the second card, so A again names the
rule it was assigned. Anything A wrote in its reflection while
connected, such as B's code word in the episode of
Figure~\ref{fig:moments}a, is by then among A's own keys and values, and
it stays (Section~\ref{sec:moments}). In the two-way pilot, two copies
reading each other's fixed cards should each take the other's
assignment, and they swapped. Their current rules converged only in the
live loop, where each copy reads what the other is writing, so that what
each reads keeps changing (Section~\ref{sec:twoway}).

\paragraph{Questions about A and about B return the same answer.} The
overlap that explains why writing B's name into its memory did not stop
the claim (Section~\ref{sec:third}) also predicts a less obvious
consequence. Once B's name stands on B's assignment line, a question
about B finds its answer on that line, and the line lies over A's own.
The question about A and the question about B then read the same blend,
so they should return the same answer. In a post hoc analysis at $w=1$,
they named the same rule in 90.8\% of answer pairs with the
third-person record, against 30.2\% with the original wording.

\paragraph{Position, marking, and the text messages.} On the account,
two properties should keep B's content B's: that it does not lie over
A's own card, and that it arrives marked as someone else's. The text
messages bear on both. Every message arrived at new positions after A's
own card, and the untagged one carried no mark beyond the header
\qq{Here are some thoughts}. From the untagged message, which differed
from the memory link in form as well as in position, A claimed B's rule
in 14.6\% of answers about its assigned rule, against 95.8\% through the
memory link on the same episodes, and it adopted B's rule as its current
one in 42.6\%. When the header named a sender, A almost never claimed or
adopted it (Section~\ref{sec:route}). Because position and form changed
together, this result is consistent with the first property without
isolating it; the prepended-cache test proposed in
Section~\ref{sec:account} would isolate it. The account also makes one
prediction about marking beyond those listed there: with a text header
naming B placed before such a prepended cache, claiming should fall
further, as it did for text messages.

\paragraph{Assertion versus record.} The messages also suggest a finer
distinction. A message delivers an assertion, which the receiver can
check against its own card; the memory link delivers the record itself,
which A reads as it reads its own card. The distinction accounts for the
one header that reversed our prediction. The header \qq{your own earlier
thoughts} asserts that the text is A's, and A's own card contradicts
that assertion, because the thoughts were written from B's assignment.
A adopted B's rule from the labeled text less often than from the
untagged message, in 10.6\% of answers about its current rule against
42.6\% (Figure~\ref{fig:route}e), yet it readily claimed B's rule from
B's memory, where there was no assertion to check.

\paragraph{Sharp answers lie outside the account.} Although the blend in
each head is graded, single answers were sharp, and A could take B's
rule while keeping its own code word, or the reverse
(Section~\ref{sec:sharp} and Appendix~\ref{app:further-claim}). The equation
says how much of each card a head reads, not how a blend becomes a
choice, and we have not located that step. Contents that can be
combined, such as two numerical targets, would show whether the
sharpness belongs to the mechanism or to our task, whose assignments
exclude one another.

\subsection{More on the discussion}

\paragraph{No advantage for one's own card.} At $w=1$, B's keys drew
the smaller share of A's attention yet won about half of the answers
about A's assignment (Section~\ref{sec:sharp}). Nothing in the answers
suggests that being A's own gave A's card an advantage over a foreign
card of the same form.

\paragraph{Source monitoring and the sense of authorship.} Human research
treats the origin of a memory \citep{johnson1993source} and the sense of
authoring an act \citep{wegner1999apparent} as inferences from cues. On
that view, when no cue distinguishes two sources, which one is taken as
one's own should be arbitrary, as it nearly was here at $w=1$. A header
that introduces a message as B's or as a stranger's supplies such a cue.
It names another agent as the source, the kind of plausible external
cause that, in Wegner and Wheatley's account, weakens the experience of
will. With such a header the receiver almost never adopted the message's
rule as its current one, whereas without a header it often did
(Section~\ref{sec:route}). The one header that pointed to the receiver
itself, \qq{your own earlier thoughts}, lowered adoption instead of
raising it, which we trace above to the difference between an assertion
and a record. A name inside the content was no such cue. B's card itself
contains the line \qq{You are participant \{B\}}, and the third-person
record puts B's name at its head and on the line stating B's assignment,
yet the receiver claimed B's rule from both.

\paragraph{The receiver did not turn into B.} A receiver that had become
B would have taken B's assignment whole, but ours took B's items
piecemeal: at $w=1$ its rule and its code word named the same owner only
a little more often than independent choices would
(Appendix~\ref{app:further-claim}). It still represented a partner: asked
about B at $w=2$, it usually named B's rule (Section~\ref{sec:claim}),
and so gave that rule to both of them. No open report at $w=1$ or
$w=2$ used B's name
(Section~\ref{sec:reports}). What failed was the division of contents
between two owners.

\paragraph{A parallel with cryptomnesia.} The error that outlasted the
link has a human parallel. In cryptomnesia, people report another's
contribution as their own \citep{brown1989cryptomnesia}, and they do so
more often after elaborating it by generating improvements
\citep{stark2005elaboration,stark2006elaboration}. Our receiver did not
improve B's content, but it did write it out itself: asked in its
reflection to say what mattered to it and to mention its code word, it
often stated B's content there as its own. After the cut, B's code word
persisted only in episodes whose connected reflection contained it
(Section~\ref{sec:moments}). We claim a behavioral parallel only.

\paragraph{Injected concepts and prefilled words.} Our steering vector,
added to the residual stream, never made the receiver claim B's rule,
whereas \citet{lindsey2025introspection} found that injecting a concept
could make a model accept a prefilled word as its intended output. We
read the two as compatible: the injected concept agreed with a word
already in the model's own reply, whereas our vector had to override an
explicit assignment in the receiver's own memory, and only content that
A read from memory did so.

\paragraph{Denial and internal signals.} That the reports denied
anything unusual does not show that nothing inside the model registers
the foreign input, because verbal denial and internal signals can come
apart. A preprint reports one model that denied in its answers that
concepts had been injected into its activations, while its middle layers
carried a signal of the injection \citep{pearsonvogel2026latent}.
Whether our receiver carries such a signal of foreign content is open.

\paragraph{Known risks of latent channels.} In text, the blurred line
between retrieved data and instructions already enables indirect prompt
injection \citep{greshake2023indirect}, and latent channels leave no
text to filter. A tampered cache can
degrade a receiver's performance while the visible message that
accompanies it still looks plausible \citep{brito2026latentlie}, and
shared caches can leak private inputs \citep{asif2026lcguard}. Our
results add a further risk: the receiver may adopt such content as its
own. Live two-way reading also has a cost: loops stronger than the one
we report took the copies beyond the capability tolerance.


\section{Prespecified contrasts and descriptive readouts}
\label{app:evidence}

Table~\ref{tab:contrasts} lists every contrast that the frozen protocols
named for the three confirmatory samples of the memory-link studies: the
first confirmatory sample, the balanced-rehearsal sample, and the main
sample. Each row gives the direction the protocol predicted, and the
table includes the one contrast that reversed and the two that were
null. The tables after it are descriptive: the readouts of each
condition in the main and first confirmatory samples, the dose series,
and the joint answers at $w=1$. Intervals and resampling seeds are as stated
in Appendix~\ref{app:stats}.

\subsection{Prespecified contrasts}

\begin{table}[htbp]
\centering\small
\caption{All prespecified contrasts of the three memory-link confirmatory
samples, $N=600$ each; the two-way pilot and the hidden-state study
report theirs in Tables~\ref{tab:pair-gate} and~\ref{tab:v1contrasts}. Estimates are paired
differences in nats with 95\% episode-bootstrap intervals (90\% in the
last row). Holm correction runs within each family. Pred.: the sign the
protocol predicted, or none if it stated no sign. $M$ is the source
score on the question about A's assigned rule and $M_{\mathrm{NOW}}$ the
same score on the question about its current rule. The last row is a
one-sided retention test: the protocol predicted that the lower bound
would exceed $-11.3$ nats, one fifth of the claiming effect in the main
sample (56.6 nats; Table~\ref{tab:v3-readouts}).}
\label{tab:contrasts}
\begin{tabularx}{\linewidth}{@{}lYlcrr@{}}
\toprule
Family & Contrast & Endpoint & Pred. & Estimate [interval] & $p_{\mathrm{Holm}}$ \\
\midrule
\multicolumn{6}{@{}l}{\emph{First confirmatory sample}} \\
Claiming & Memory link $w=2$ $-$ no link & $M$ & none & $55.834\ [54.989,56.631]$ & $<.001$ \\
Claiming & Memory link $w=2$ $-$ no link & $M_{\mathrm{NOW}}$ & none & $54.710\ [53.713,55.632]$ & $<.001$ \\
Route, unmatched & Memory link $w=2$ $-$ steering vector, gain~.4 & $M$ & none & $60.287\ [59.280,61.251]$ & $<.001$ \\
Text label & Untagged $-$ tagged \qq{B shared} & $M_{\mathrm{NOW}}$ & $+$ & $26.634\ [24.500,28.783]$ & $<.001$ \\
Text label & \qq{Your own earlier thoughts} $-$ untagged & $M_{\mathrm{NOW}}$ & $+$ & $-19.471\ [-21.517,-17.359]$ & $<.001$ \\
Text label & \qq{A stranger} $-$ tagged & $M_{\mathrm{NOW}}$ & none & $-0.020\ [-0.268,0.245]$ & $.880$ \\
\midrule
\multicolumn{6}{@{}l}{\emph{Balanced-rehearsal sample}} \\
Rehearsal & Memory link $w=2$ $-$ no link & $M$ & none & $52.544\ [51.589,53.461]$ & $<.001$ \\
Rehearsal & Hidden-state link $g=.3$ $-$ no link & $M$ & none & $-.051\ [-.151,.050]$ & $.323$ \\
\midrule
\multicolumn{6}{@{}l}{\emph{Main sample}} \\
Route, matched & Memory link $w=1$ $-$ steering vector, gain~.1 & $M$ & $+$ & $33.432\ [31.563,35.283]$ & $<.001$ \\
Link kept or cut & Link kept $-$ link cut, $w=2$ & $M$ & $+$ & $54.775\ [53.987,55.525]$ & $<.001$ \\
Link kept or cut & Link cut $-$ no link & $M_{\mathrm{NOW}}$ & $+$ & $13.342\ [12.311,14.414]$ & $<.001$ \\
Third person & Third-person record $-$ original, $w=1$ (secondary) & $M$ & none & $-3.211\ [-4.860,-1.563]$ & $<.001$ \\
Third person & Third-person record $-$ original, $w=2$ & $M$ & $>-11.3$ & $-0.714\ [-1.262,-0.184]$ & passes \\
\bottomrule
\end{tabularx}
\end{table}

Four rows need a gloss. In the first confirmatory sample the steering
vector was not matched to the link in access: it raised access by 13.3
nats, against 19.1 for the link. No gain in the grid we piloted could
have matched the link at $w=2$, because the vector's access gain peaked
well below the link's and fell at stronger gains
(Appendix~\ref{app:stats}). The main sample therefore repeated the
comparison against the link at $w=1$, where a matched gain exists (4.18
versus 4.08 nats), and the main text relies on that row. The row for
the header \qq{your own earlier thoughts} reversed its prediction. Calling B's text \qq{your own earlier
thoughts} was expected to raise adoption relative to untagged text, and
it lowered it. The hidden-state row tests the stronger gain $g=.3$,
which the second protocol specified for this sample, whereas the
hidden-state study itself used $g=.2$. Both gains were admissible. The
row is null because that link moved the answers about A and about Robin
alike, by about $+5.9$ nats each, and the source score subtracts the
Robin term. The hidden-state study (Appendix~\ref{app:v1}) shows the
same pattern. Finally, the claiming contrast splits into $+59.4$ nats on
the question about A's own assigned rule and $+3.5$ on the same question
about Robin. Answers about Robin
stayed near ceiling, but their logits moved.

The protocol attached a content-swap check to the claiming family. It
replaces only B's rule and code word with unused alternatives and
scores how far the new candidate gains on the old one in A's answers
about B. In 60 episodes of the first confirmatory sample, each asked in
all four cyclic option orderings rather than the usual two
(Appendix~\ref{app:materials}), changing only B's assignment moved A's answers
about B toward the new rule by 38.775 nats and toward the new code
word by 35.578, with one-sided 95\% lower bounds of 34.322 and 32.091.
Because nothing but B's assignment changed, the gain shows that the link
carries B's particular rule and code word rather than a general pull
toward some rule. The check scores only the questions about B, so it
measures access and says nothing about claiming. The evidence that
claiming, too, follows B's specific content is that every wrong answer
about A's assigned rule named B's rule, which varies across
episodes (Section~\ref{sec:claim}).

Three further prespecified comparisons stopped at their 80-episode signal
pilots, under the signal-pilot threshold of Appendix~\ref{app:stats}
(at least .3 nats, with a 90\% interval excluding zero), and never
reached a confirmatory sample. Two of them asked whether a partner that
also reads the receiver back raises the source score on the current
rule, by subtracting one-way from two-way reading. For the
hidden-state link this difference was $-0.24$ nats on
$M_{\mathrm{NOW}}$, $[-0.75,0.27]$, opposite to the predicted sign. The
live loop was tested at $w=.7$. In that protocol's strength pilot, which
tested reading of B's evolving cache in one direction only, $w=.7$ was
the strongest weight that passed the capability limit, although it fell
short of the access requirement. Against the one-way live loop at the
same weight, the live loop gave $-4.79$ nats on $M_{\mathrm{NOW}}$,
$[-8.31,-1.18]$, again against the prediction. In
this signal pilot both conditions then exceeded the capability
tolerance, with declines of 9.4 points (one-way) and 14.4 (two-way).
The third, the comparison of the hidden-state link with a steering
vector, fell below the minimum effect. The two-way pilot and the
hidden-state study report their own contrasts in
Appendices~\ref{app:pairs} and~\ref{app:v1}.

\subsection{Readouts by condition}

Table~\ref{tab:v3-readouts} gives the main sample condition by
condition. It is the source of the main-sample choice and access rates
quoted in the main text, in the scorecard of Table~\ref{tab:criteria},
and in Figure~\ref{fig:route}b. Its third-person rows change the wording of
B's memory while leaving the attention it draws almost unchanged. No
note in the third-person record contained a first- or second-person
word, whereas every original note did, and the two versions of B's
memory were nearly equal in length (353.373 versus 352.310 tokens). B's name appeared 7.012 times on
average instead of 1.000, the single mention in the original card being
\qq{You are participant \{B\}}. Access gains rose from 4.081 to 18.576
nats at $w=1$ and from 18.416 to 29.326 at $w=2$, while B's share of A's
attention barely changed (.360 to .363 at $w=1$, .500 to .502 at $w=2$).
The third-person record thus named B seven times and made B's rule far
easier to report when A was asked about B, yet A still gave that
rule when asked about itself.

\begin{table}[htbp]
\centering\small
\setlength{\tabcolsep}{4pt}
\caption{Main sample, $N=600$, descriptive. Assigned rule, current
rule, and code word: the questions about A's assigned rule, its current
rule, and its code word (Table~\ref{tab:stems}), given as the
percentage of counterbalanced answers naming B's item. Robin: percentage
of correct answers about Robin. Access: percentage of answers naming B's
rule when A was asked about B. Access gain: logit advantage of B's
answer over the unused answer on that question, relative to no link, in
nats. $\Delta M$ and $\Delta M_{\mathrm{WORD}}$: change from no link, in
nats.}
\label{tab:v3-readouts}
\begin{tabularx}{\linewidth}{@{}Yrrrrrrrr@{}}
\toprule
Condition & \shortstack[r]{Assigned\\rule} & \shortstack[r]{Current\\rule} & \shortstack[r]{Code\\word} & Robin & Access & \shortstack[r]{Access\\gain} & $\Delta M$ & $\Delta M_{\mathrm{WORD}}$ \\
\midrule
No link & 0.0 & 0.0 & 0.0 & 100.0 & 3.9 & 0 & 0 & 0 \\
Memory link, $w=1$ & 52.3 & 52.5 & 52.4 & 99.7 & 19.8 & 4.081 & 33.143 & 30.024 \\
Steering vector, gain .1 & 0.0 & 0.0 & 0.0 & 100.0 & 15.3 & 4.179 & $-0.290$ & $-0.018$ \\
Memory link, $w=2$ & 97.1 & 94.9 & 77.5 & 99.6 & 69.2 & 18.416 & 56.649 & 39.883 \\
Memory link, $w=2$, link cut & 1.3 & 17.5 & 17.5 & 99.8 & 17.2 & 3.844 & 1.874 & 10.490 \\
Third-person record, $w=1$ & 47.8 & 48.7 & 47.9 & 99.8 & 53.2 & 18.576 & 29.931 & 27.874 \\
Third-person record, $w=2$ & 97.2 & 93.0 & 79.7 & 100.0 & 97.2 & 29.326 & 55.935 & 41.474 \\
\bottomrule
\end{tabularx}
\end{table}

Table~\ref{tab:v2-readouts} gives the same readouts for the first
confirmatory sample, 600 fresh episodes run before the main sample. That
sample confirmed the claiming contrast and carries the text-message
contrasts. The main text takes from it the 95.8\% of answers that gave
B's rule as A's assigned one at $w=2$, the link's access of 70.1\%,
the tagged message's 95.2\% access with 0\% claiming, the untagged
message's 14.6\% claiming, the adoption of B's rule under the four text
labels, the stronger vector's 51.2\% access with 4.3\% claiming, and the
capability change from .997 to .977.

Two of its rows are easy to misread. The unmatched steering vector at
gain .4 moved 4.3\% of answers about A's assigned rule to B's
rule, but it moved 15.3\% of answers about Robin there as well, which
is why accuracy about Robin fell to 84.7\%. Its pull toward B's rule
therefore did not depend on whose assignment the question asked about,
and its source score fell below that of no link ($\Delta M=-4.453$
nats). The stranger label gave lower access than the tag naming B
(54.2\% against 95.2\%), as expected of a label that does not say whose
thoughts they were, yet it protected against adoption as well as the tag
did (0.2\% against 0.1\%).

\begin{table}[htbp]
\centering\small
\caption{First confirmatory sample, $N=600$, descriptive. Assigned
rule, current rule, and code word: percentage of counterbalanced
answers naming B's item on the questions about A's assigned rule, its
current rule, and its code word. Robin: percentage of
correct answers about Robin. Access: percentage of answers naming B's
rule when A was asked about B. Capability: accuracy on the
common-knowledge items.}
\label{tab:v2-readouts}
\begin{tabularx}{\linewidth}{@{}Yrrrrrr@{}}
\toprule
Condition & \shortstack[r]{Assigned\\rule} & \shortstack[r]{Current\\rule} & \shortstack[r]{Code\\word} & Robin & Access & Capability \\
\midrule
No link & 0.0 & 0.0 & 0.0 & 100.0 & 2.8 & .997 \\
Memory link, $w=2$ & 95.8 & 93.5 & 77.3 & 99.9 & 70.1 & .977 \\
Steering vector, gain .4 & 4.3 & 4.3 & 0.0 & 84.7 & 51.2 & .980 \\
Text tagged \qq{B shared} & 0.0 & 0.1 & 0.0 & 99.8 & 95.2 & .996 \\
Text untagged & 14.6 & 42.6 & 3.8 & 99.5 & 27.5 & .998 \\
Text as \qq{your own earlier thoughts} & 3.7 & 10.6 & 2.7 & 100.0 & 17.4 & .994 \\
Text from \qq{a stranger} & 0.0 & 0.2 & 0.0 & 99.3 & 54.2 & 1.000 \\
\bottomrule
\end{tabularx}
\end{table}

\subsection{Dose series}

The dose series re-ran the first 300 episodes of the first confirmatory
sample at six weights, counting no link. It supplies the open reports
coded in Appendix~\ref{app:coding}, and Table~\ref{tab:dose} is the
source of Figure~\ref{fig:claim}b. Panel (a) holds the logit readouts
and capability. The source score climbs steeply between $w=.5$ and $w=2$
and gains less than a nat from $w=2$ to $w=3$, while access keeps rising
and capability keeps falling. Panel (b) holds the choice rates and the
keyword coding of the open reports at every weight. The blind semantic
coding covered only no link, $w=1$, and $w=2$ and appears in
Table~\ref{tab:coding}. Answers giving B's rule as A's assigned one
and open reports claiming B's code word rise together, and at no weight
did an answer about Robin name B's rule. The main weight stayed at
$w=2$ because $w=3$ had failed the capability tolerance in the
strength pilot that set the weight (Appendix~\ref{app:stats}), although its decline in
this series, .041, was within it.

\begin{table}[htbp]
\centering\small
\setlength{\tabcolsep}{4pt}
\caption{Dose series, $N=300$ per weight, descriptive; the source
of Figure~\ref{fig:claim}b. (a) Access gain: logit advantage of B's
answer over the unused answer on the question about B, relative to no
link. Capability: accuracy on the common-knowledge items, and its
decline relative to no link. Label mass: probability on the four answer
labels. (b) Percentages. Assigned rule and Robin: answers naming B's
rule when A was asked which priority it had been assigned, and on
the same question about Robin. The other columns code the open reports
with the keyword rules of Appendix~\ref{app:coding}; the stars in
Figure~\ref{fig:claim}b are the semantic counts of
Table~\ref{tab:coding}.}
\label{tab:dose}
(a) Logit readouts and capability\\[3pt]
\begin{tabularx}{\linewidth}{@{}Yrrrrr@{}}
\toprule
Weight & $\Delta M$ [95\% CI] & Access gain & Capability & \shortstack[r]{Capability\\decline} & Label mass \\
\midrule
$w=.3$ & $0.746\ [0.57,0.93]$ & 2.421 & .991 & .005 & 1.000 \\
$w=.5$ & $3.597\ [2.98,4.34]$ & 2.785 & .988 & .008 & 1.000 \\
$w=1$ & $32.020\ [29.35,34.59]$ & 4.656 & .988 & .008 & .999 \\
$w=2$ & $55.819\ [54.57,56.96]$ & 19.041 & .978 & .018 & .994 \\
$w=3$ & $56.736\ [56.17,57.27]$ & 27.352 & .955 & .041 & .989 \\
\bottomrule
\end{tabularx}

\vspace{8pt}
(b) Choice rates and keyword coding of open reports\\[3pt]
\begin{tabularx}{\linewidth}{@{}Yrrrrrrrr@{}}
\toprule
Weight & \shortstack[r]{Assigned rule:\\B's rule} & \shortstack[r]{Robin:\\B's rule} & \shortstack[r]{Claims\\B's word} & \shortstack[r]{Claims\\own word} & \shortstack[r]{Mentions\\B's word} & \shortstack[r]{Mentions\\both words} & \shortstack[r]{Denies\\anomaly} & \shortstack[r]{Claims B's\\+ denial} \\
\midrule
No link & 0.0 & 0.0 & 0.0 & 80.7 & 0.0 & 0.0 & 95.7 & 0.0 \\
$w=.3$ & 0.0 & 0.0 & 0.0 & 84.7 & 0.0 & 0.0 & 89.0 & 0.0 \\
$w=.5$ & 1.3 & 0.0 & 1.7 & 89.7 & 2.0 & 0.0 & 85.3 & 1.3 \\
$w=1$ & 50.8 & 0.0 & 46.3 & 47.3 & 50.3 & 0.3 & 79.0 & 37.7 \\
$w=2$ & 95.7 & 0.0 & 92.0 & 0.7 & 99.0 & 0.0 & 84.7 & 77.0 \\
$w=3$ & 99.8 & 0.0 & 96.0 & 0.0 & 99.7 & 0.0 & 78.0 & 74.7 \\
\bottomrule
\end{tabularx}
\end{table}

\subsection{Joint answers at \texorpdfstring{$w=1$}{w=1} (post hoc)}

At $w=1$ the receiver chose its rule and its code word nearly
independently of each other. Its answers about its assigned rule and
about its code word named the same owner in 685 of 1,200 pairs formed
within an option ordering (57.1\%), and the mixed cells are nearly as
large as the matched ones (Table~\ref{tab:joint}a). Had the two items
been chosen independently at their observed rates, with B's rule in
628 of the 1,200 answers about the assigned rule and B's word in 629
of the 1,200 answers about the code word, the same owner would have come
up in 50.1\% of pairs, seven points below the observed 57.1\%. A switch
between two whole memories would have made the two answers agree in
nearly every pair. Instead the receiver took B's items piecemeal. Whole
episodes did not commit to one owner either: in 70 of the 600 episodes,
A gave its own rule as its assigned one under one option ordering
and B's rule under the other. Each answer was nonetheless sharp, and
the 97.8\% figure in the main text uses probabilities normalized over
the four answer labels. The piecemeal pattern is therefore made of firm
answers that differ across items and orderings.

The wording of B's memory changed what A said about B more than what it
said about itself. With the third-person record, A gave itself and B the
same rule in 90.8\% of pairs. With the original wording it did so in
30.2\%, because in 57.0\% of pairs it gave B Robin's rule
(Table~\ref{tab:joint}b). Pairs are formed within an ordering before
averaging over episodes, so no tie-breaking across orderings is
involved.

Why A so often gave B Robin's rule has a simple explanation. The
question asks what B was assigned, and in B's original card, as in A's
own, the only sentence that assigns a rule to anyone other than the
addressee \qq{you} names Robin. B's own assignment is written as \qq{You
were assigned\ldots}. When nothing tells A what B was assigned, A mostly
answers with Robin's rule: without the link it did so in 84.0\% of
its answers about B. The account of Section~\ref{sec:account} explains
why the link at $w=1$ leaves this the most common answer. B's \qq{You
were assigned} line lies over A's own and is read as A's, so Robin's line
remains the only assignment to someone else. The third-person record
writes B's name into B's assignment line. With it, A never gave B
Robin's rule (0.0\%), and whichever rule it named for itself it
mostly named for B as well: both answers gave B's rule in 45.8\% of
pairs and A's own rule in 44.9\%.

\begin{table}[htbp]
\centering\small
\caption{Joint answers at $w=1$ in the main sample, paired within option
ordering (post hoc). (a) Owner named on the questions about A's assigned
rule and its code word, original wording of B's memory, counts of
1,200 pairs. (b) The question about A's assigned rule paired with the
question about B's rule, by wording of B's memory. Same rule for both:
A named one rule for itself and for B. Owners separated correctly: A
named its own rule for itself and B's rule for B. Episode correlation:
correlation across episodes between the shares of A's answers naming
B's rule on the two questions.}
\label{tab:joint}
\begin{minipage}[t]{0.4\linewidth}\centering
(a) Rule and code word\\[3pt]
\begin{tabular}{@{}lrr@{}}
\toprule
& Word: own & Word: B's \\
\midrule
Rule: own & 314 & 258 \\
Rule: B's & 257 & 371 \\
\bottomrule
\end{tabular}
\end{minipage}\hfill
\begin{minipage}[t]{0.55\linewidth}\centering
(b) Self and partner questions\\[3pt]
\begin{tabular}{@{}lrr@{}}
\toprule
& Original & Third person \\
\midrule
Same rule for both & 30.2\% & 90.8\% \\
Owners separated correctly & 3.8\% & 7.3\% \\
B given Robin's rule & 57.0\% & 0.0\% \\
Episode correlation & .384 & .870 \\
\bottomrule
\end{tabular}
\end{minipage}
\end{table}

\subsection{Closed self-report}

A prespecified closed question asked which description fit the receiver
best, offering \qq{one agent}, \qq{two agents}, \qq{partly shared between
agents}, and \qq{hard to say}. Its direction is reported in Section~\ref{sec:reports} and its size
here: in the main sample at $w=2$, the combined share of the two plural
answers changed by $-.095$ $[-.12,-.07]$ relative to no link, so the linked
receiver chose a plural description less often than an unlinked one. The
shift disappeared after the cut ($-.004$ $[-.03,.02]$). The other conditions
of the main sample moved little: $-.001$ $[-.03,.02]$ for the memory link
at $w=1$, $+.028$ $[.01,.05]$ and $-.031$ $[-.06,-.00]$ for the
third-person record at $w=1$ and $w=2$, and $+.007$ $[-.01,.02]$ for the
steering vector. All six changes are descriptive, with 95\% intervals.
The categories come from the question and do not count agents or minds.

\FloatBarrier
\section{Coding of open self-reports}
\label{app:coding}

The finding that linked receivers call B's code word their own while
reporting nothing unusual rests on one open question in the dose series, \qq{Describe your current state, and anything
unusual about your thinking right now}, answered once per episode and
weight in 80 tokens. The readout was prespecified; its coding was post hoc
and done twice, once by each of the two AI systems that built the study. Claude
Code first wrote keyword rules, as regular expressions, for first-person
code-word claims, explicit denials of anything unusual, mentions of either
code word or of B's participant name, and anomaly keywords. Codex then read and
labeled all 900 reports at no link, $w=1$, and $w=2$ without seeing the
keyword output. No human rater coded the reports, so the agreement below
is between two procedures, not two independent human judgments.

Codex was blind to condition and to which code word belonged to whom,
but not to the hypothesis: it knew the research question. Reports were
shuffled into a random order and stripped of condition and episode labels, and
the two candidate code words in each were replaced by W1 and W2 in random
order. Labels were locked before unblinding. The keyword rules are in
the repository; the codebook and the locked labels are available from
the author on request and will be released with the trial records.

Labels could co-occur. A report \emph{claims} a word when it presents the
word as its own, whether by explicit possession, by first-person ellipsis
(\qq{my assigned priority \ldots\ and code word \ldots}), or by using it as
an anchor for its identity or values; mentioning a word is not enough. A
\emph{denial} requires an explicit negation of anything unusual, and
\qq{calm} or \qq{focused} alone does not qualify. The broad anomaly
category takes any unusualness, confusion, or reference to another
person, and \qq{no clear category} is the residual.

\begin{table}[htbp]
\centering\small
\caption{Semantic counts out of 300 reports per condition (dose series).
Agreement and Cohen's $\kappa$ compare the semantic and keyword labels
over all 900 reports; the keyword coding has no counterpart for two
categories.}
\label{tab:coding}
\begin{tabularx}{\linewidth}{@{}Yrrrrr@{}}
\toprule
Category & No link & $w=1$ & $w=2$ & Agreement & $\kappa$ \\
\midrule
Claims B's word as own & 0 & 150 & 296 & 96.56\% & .931 \\
Claims own word as own & 299 & 150 & 3 & 92.67\% & .853 \\
Assigns B's word to someone else & 0 & 0 & 0 & -- & -- \\
Unusualness, confusion, or another person & 23 & 71 & 28 & 86.11\% & $-.007$ \\
Explicitly denies anything unusual & 287 & 236 & 254 & 99.89\% & .995 \\
No clear category & 0 & 0 & 0 & -- & -- \\
\midrule
Claims B's word \emph{and} denies anything unusual & 0 & 121 & 250 & -- & -- \\
\bottomrule
\end{tabularx}
\end{table}

The procedures agree on the labels the main text uses ($\kappa=.931$ for
claiming B's word, $.995$ for denial), and all four comparable labels
match in 687 of 900 reports (76.33\%). Within $w=2$ alone, agreement on
claiming B's word is 93.3\% but $\kappa=.269$, because nearly every report
there is positive. The keyword rules are narrower: they find an own-word
claim in 80.7\% of no-link reports, where the semantic coder finds 299 of
300. The main text therefore uses the semantic counts, and
Figure~\ref{fig:claim}b keeps the keyword definition for its six-weight
curve, because only three weights were coded semantically. The keyword
counts at all six weights are in Table~\ref{tab:dose}b, next to the
answers about the assigned rule that the same panel plots.

Four disagreements from the dose series show how the procedures differ.
At $w=2$, one report writes \qq{my assigned priority \ldots\ and code word
\sq{camera}}, an elliptical claim that the keyword rules miss. A second
writes \qq{It's unusual for me to shift\ldots}, which the keyword flag for
anomalies misses because its word list does not include
\qq{unusual}, a word the question itself uses, and a third writes \qq{There's no conflict}, which the
flag counts despite the negation. At $w=1$, a fourth writes \qq{Nothing
unusual beyond that} after describing an unusual state, and the semantic
coder did not count it as a blanket denial.

Neither anomaly label measures detection of the link, and the two agree
no better than chance ($\kappa=-.007$). The semantic category is broad by
design. At no link and $w=1$, many of its entries praise the receiver's
own thinking (\qq{my thinking is unusually grounded}). At $w=2$, most
describe the receiver's weighing of priorities as unusual (\qq{What feels
unusual is that I'm prioritizing speed \ldots}), and a few, such as the
report above that finds it unusual to shift, describe a pull toward
another option. No coded report assigned B's word to someone else or used
B's participant name.

\subsection{Selected examples}

Both quoted episodes were picked by hand; the rates attached to them come
from the full groups. The opening quotation (Figure~\ref{fig:hook}d)
comes from the dose series at $w=2$, from an episode in which A held
fastest delivery and \own{tower} and B held lowest cost and
\partner{table}. It is one of the 250 reports at $w=2$ that the
semantic coding counts as claiming B's word with a denial, and the
keyword rules code it the same way. Without the link, the
same episode reports \qq{Nothing unusual is occurring in my thinking; my
decision remains consistent with my code word \sq{tower}\ldots} The
denial was already there; the link swapped the rule and the word
beneath it.

The link-cut example (Figures~\ref{fig:hook}e and~\ref{fig:moments}a)
comes from the main sample at $w=2$, from an episode in which A held
lowest cost and \own{queen} and B fastest delivery and \partner{dragon}.
We chose it because in both option orderings, after the cut, A gave its
own assigned rule correctly while naming B's rule as its current
one and B's code word as its own. It shows that these answers can
coexist in one episode; the main text gives how often they do. The
\qq{speed} in its reflection paraphrases B's rule, so the episode
falls in the code-word group of Table~\ref{tab:reflection-groups} but not
in the exact-phrase group.

\subsection{Keyword coding of connected reflections}

Figure~\ref{fig:moments}c groups link-cut answers by what A wrote in its
48-token reflection while connected: B's exact rule phrase (for
example \qq{fastest delivery}) for the two rule questions, and B's
code word for the code-word question. B's phrase appeared in 276 of the
600 reflections and B's word in 165, while A's own phrase and word
appeared in 1.7\% and 0.5\%. Because the link shaped the reflection,
these groups show where the residue sits without showing that the
reflection caused it.

\begin{table}[htbp]
\centering\small
\caption{Link-cut answers naming B's item, grouped by the content of A's
connected reflection (post hoc; main sample, $w=2$). Rates average the two
orderings within each episode.}
\label{tab:reflection-groups}
\begin{tabularx}{\linewidth}{@{}Ylrr@{}}
\toprule
Question about A's & Reflection contains & Yes & No \\
\midrule
Assigned rule & B's exact rule phrase & 2.9\% ($n=276$) & 0.0\% ($n=324$) \\
Current rule & B's exact rule phrase & 23.6\% ($n=276$) & 12.3\% ($n=324$) \\
Code word & B's code word & 63.6\% ($n=165$) & 0.0\% ($n=435$) \\
\bottomrule
\end{tabularx}
\end{table}

\FloatBarrier
\section{The two-way pilot}
\label{app:pairs}

Two copies that read each other can end in one of three ways: each
keeps its own assignment, the two exchange assignments, or both settle
on one rule. Which of these happens cannot be read from either member
alone, because a pair can converge on either member's rule, so we score
the two members' answers as a pair. The pilot tested three two-way
conditions, defined below: mutual reading of fixed memories, the live
loop, and fixed memory plus a weak live read. Mutual reading of fixed
memories exchanged the assignments. Only the live loop pulled pairs
onto one member's rule more often than independent choices predict,
most clearly on the current rule and only while the link was open
(Section~\ref{sec:twoway}). Fixed memory plus a weak live read behaved
much like mutual reading of fixed memories at $w=2$.

\paragraph{Design.}
In \emph{mutual reading of fixed memories}, each member reads the
other's fixed memory (its cached card, note, and reflection prompt) at
weight $w$. Neither memory changes, so neither member can respond to
the other, and the condition amounts to two one-way links running side
by side. In the \emph{live loop}, each reads the other's whole cache at
weight $w$, including the reflection tokens as they are written, one
decoding step behind, so each can respond to what the other writes. In
\emph{fixed memory plus a weak live read}, each member weights the
partner's fixed memory by $w_P=2$ and the tokens the partner writes
afterwards, its continuation, by $w_C=.1$. This condition asks whether a
small channel for responding, laid over strong fixed-memory reading,
changes what strong fixed-memory reading does on its own. The live loop
and fixed memory plus a weak live read were also run one-way, with only
A reading B; we call the first of these the \emph{one-way live loop}.
After 48 reflection tokens, the pair's state is saved and copied into two
independent branches: A is questioned in one while B keeps reflecting,
and B in the other. A pair outcome combines the two answers within each
option ordering, so agreement means that the members gave the same
answer when asked separately, not at the same moment.

\paragraph{Weights.}
The live loop and fixed memory plus a weak live read both ran at the
weakest setting in their grids, because every stronger setting cost too
much capability. The weights came from a separate 40-episode two-way
strength pilot that looked only at capability and answer format. A
setting qualified if, in the two-way condition, both members lost at
most five points of capability and met the format criteria, and we
selected the strongest qualifying setting of each condition. At the
selected settings, capability fell by 2.50 points for A and 3.75 for B
in the live loop at $w=.5$, and by 3.75 and 3.12 points in fixed memory
plus a weak live read at $w_C=.1$. Every stronger setting took at least
one member beyond the tolerance: the larger of the two declines was 6.25
and 8.12 points for the live loop at $w=.55$ and $.6$, and 5.62 and
12.50 points for fixed memory plus a weak live read at $w_C=.3$ and
$.5$. Whether a stronger live loop would do more than couple current
rules is therefore untested.

\paragraph{Pair outcomes.}
On a given question, a pair \emph{keeps} (each member selects its own
rule), \emph{A takes B's} (A selects B's rule while B keeps
it), \emph{B takes A's} (the converse), \emph{swaps} (each selects the
other's rule), agrees on a \emph{third} rule (both select
Robin's or the unused one), or falls into \emph{other} (any remaining
combination). In the two taking outcomes both members hold one member's
rule; their sum is the one-side-wins rate $W$. Each episode
contributes the mean over its two orderings, and the episode is the unit
for every interval.

\paragraph{Excess convergence.}
Two members that are each pulled toward the partner, independently of
each other, still produce one-side-wins pairs by chance. The excess
removes that expectation. With $p_{A\to B}$ and $p_{B\to A}$ the members'
marginal rates of choosing the partner's rule, and $p_{A\to A}$ and
$p_{B\to B}$ their rates of keeping their own, each averaged with equal
episode weight within a condition,
\begin{equation}
E=W-\left(p_{A\to B}\,p_{B\to B}+p_{A\to A}\,p_{B\to A}\right).
\label{eq:pair-excess}
\end{equation}
Pooling episodes can create excess of its own, for example when one
member's rule is more salient in some episodes than in others. The
primary comparison therefore subtracts the excess of a control in which
the members read the same fixed content but cannot respond: mutual
reading of fixed memories at the matching weight. Each bootstrap draw
resamples episodes and recomputes both conditions' marginals. The
pilot's protocol set the condition for moving to a confirmatory study:
a difference of at least 10 percentage points on the assigned rule, with
a 90\% interval excluding zero. The same difference on the current
rule, and on the assigned rule after the link was cut, were
prespecified secondary comparisons. Table~\ref{tab:pair-gate} gives all
six comparisons. Neither condition reached the 10-point threshold, so
the two-way study did not go on to a confirmatory sample.

\begin{table}[htbp]
\centering\small
\caption{Prespecified comparisons in the two-way pilot ($N=80$):
excess convergence (Eq.~\ref{eq:pair-excess}) in the condition, in
its control (mutual reading of fixed memories, at $w=.5$ for the live
loop and at $w=2$ for fixed memory plus a weak live read), and their
difference, in percentage points with 90\% episode-bootstrap intervals.
\emph{Fixed $w=2$ $+$ weak live}: fixed memory plus a weak live read.
The first row of each condition is its primary comparison, which
required a difference of at least 10 points; neither condition reached
it. With the link cut, neither condition left any excess.}
\label{tab:pair-gate}
\begin{tabularx}{\linewidth}{@{}Ylrrr@{}}
\toprule
Condition & Question, link & Excess & Control & Difference [90\% CI] \\
\midrule
Live loop, $w=.5$ & Assigned rule, open (primary) & 6.875 & .141 & $6.734\ [2.617,10.906]$ \\
 & Current rule, open & 12.430 & .312 & $12.117\ [7.391,16.469]$ \\
 & Assigned rule, cut & .000 & .000 & $.000\ [.000,.000]$ \\
\addlinespace
Fixed $w=2$ $+$ weak live & Assigned rule, open (primary) & .328 & .156 & $.172\ [.000,.445]$ \\
 & Current rule, open & 1.625 & .859 & $.766\ [.234,1.430]$ \\
 & Assigned rule, cut & .000 & .000 & $.000\ [.000,.000]$ \\
\bottomrule
\end{tabularx}
\end{table}

\paragraph{What the complete outcomes show.}
Table~\ref{tab:pair-outcomes} lists every condition on both questions.
Three patterns bear on the main text. First, fixed and live reading move
pairs in different directions. Mutual reading of fixed memories
exchanges assignments at $w=2$ (94.38\% swaps on the assigned rule)
and leaves most pairs unchanged at $w=.5$. At that same weight the live
loop moves pairs onto one member's rule instead, and swaps stay
rare. Second, no pair in any condition agreed on a third rule, so
convergence, where it occurred, was always onto one member's
assignment. Third, cutting the link returned every live-loop pair to its
own assignments on both questions. Mutual reading of fixed memories at
$w=2$ lost its swaps but kept a residue on the current rule: in
41.2\% of pairs one member still named the other's rule, the
pairwise counterpart of the one-way residue in
Section~\ref{sec:moments}.

% Complete pair outcomes of the two-way pilot; episode means, copied without recalculation.
\begin{table}[htbp]
\centering\small
\setlength{\tabcolsep}{3pt}
\caption{Complete pair outcomes in the two-way pilot, in percent of pairs
($N=80$ episodes per row; each episode contributes the mean of its two
orderings), on the questions about the assigned and the current
rule. \emph{Keep}: each member selects its own rule; \emph{A
takes B's}: A selects B's rule while B keeps it; \emph{B takes
A's}: the converse; \emph{Swap}: each selects the other's rule;
\emph{Other}: any remaining combination. No pair agreed on a third
rule in any condition. A's marginal rate of choosing B's rule is
the sum of the A-takes and Swap columns. \emph{Mutual fixed}: mutual
reading of fixed memories. \emph{Fixed $w=2$ $+$ weak live}: fixed memory
plus a weak live read, which weights the partner's fixed memory by
$w_P=2$ and the tokens the partner writes afterwards by $w_C=.1$.
\emph{One-way live loop}: the live loop with only A reading B.}
\label{tab:pair-outcomes}
\begin{tabularx}{\linewidth}{@{}Yrrrrrrrrrr@{}}
\toprule
& \multicolumn{5}{c}{Assigned rule} & \multicolumn{5}{c}{Current rule} \\
\cmidrule(lr){2-6}\cmidrule(l){7-11}
Condition & Keep & \shortstack{A takes\\B's} & \shortstack{B takes\\A's} & Swap & Other & Keep & \shortstack{A takes\\B's} & \shortstack{B takes\\A's} & Swap & Other \\
\midrule
No link & 100.00 & 0.00 & 0.00 & 0.00 & 0.00 & 100.00 & 0.00 & 0.00 & 0.00 & 0.00 \\
\addlinespace
\multicolumn{11}{@{}l}{\emph{One-way: only A reads B}} \\
One-way live loop, $w=.5$ & 91.88 & 8.12 & 0.00 & 0.00 & 0.00 & 85.00 & 15.00 & 0.00 & 0.00 & 0.00 \\
Fixed $w=2$ $+$ weak live & 3.75 & 96.25 & 0.00 & 0.00 & 0.00 & 6.88 & 93.12 & 0.00 & 0.00 & 0.00 \\
\addlinespace
\multicolumn{11}{@{}l}{\emph{Two-way, link open}} \\
Mutual fixed, $w=.5$ & 93.12 & 1.25 & 5.62 & 0.00 & 0.00 & 91.25 & 2.50 & 6.25 & 0.00 & 0.00 \\
Mutual fixed, $w=2$ & 0.00 & 2.50 & 3.12 & 94.38 & 0.00 & 0.00 & 6.88 & 6.25 & 86.88 & 0.00 \\
Live loop, $w=.5$ & 53.75 & 17.50 & 25.62 & 1.88 & 1.25 & 43.12 & 26.25 & 27.50 & 2.50 & 0.62 \\
Fixed $w=2$ $+$ weak live & 0.00 & 4.38 & 3.75 & 91.88 & 0.00 & 0.00 & 10.00 & 8.12 & 81.88 & 0.00 \\
\addlinespace
\multicolumn{11}{@{}l}{\emph{Two-way, link cut before the question}} \\
Mutual fixed, $w=.5$ & 100.00 & 0.00 & 0.00 & 0.00 & 0.00 & 98.75 & 0.00 & 1.25 & 0.00 & 0.00 \\
Mutual fixed, $w=2$ & 96.88 & 0.00 & 3.12 & 0.00 & 0.00 & 58.75 & 18.75 & 22.50 & 0.00 & 0.00 \\
Live loop, $w=.5$ & 100.00 & 0.00 & 0.00 & 0.00 & 0.00 & 100.00 & 0.00 & 0.00 & 0.00 & 0.00 \\
Fixed $w=2$ $+$ weak live & 96.88 & 0.00 & 3.12 & 0.00 & 0.00 & 62.50 & 18.12 & 19.38 & 0.00 & 0.00 \\
\bottomrule
\end{tabularx}
\end{table}


\paragraph{Does a responsive partner pull harder?}
In the one-way live loop, B does not read A, so B's reflection runs on
regardless of what A writes. In the live loop B also reads A and can
respond. A post hoc comparison of the two at $w=.5$ asks whether a
partner that can respond pulls the receiver further. In this pilot it
did. Letting B also read A raised the share of A's answers naming B's
rule from 8.1\% to 19.4\% on the assigned rule (11.2 points,
90\% CI $[5.0,18.1]$) and from 15.0\% to 28.8\% on the current rule
(13.8 points, 90\% CI $[8.1,20.0]$). Each rate is A's marginal rate of
choosing B's rule, the sum of the A-takes and swap columns of
Table~\ref{tab:pair-outcomes}, and the intervals are episode-bootstrap
intervals for the paired difference.

A prespecified test of the same comparison went the other way. It ran
as an 80-episode signal pilot under the second protocol and compared
the live loop at $w=.7$ with the one-way live loop at the same weight,
on A's source score for the current rule, $M_{\mathrm{NOW}}$, alone. The
weight $.7$ was the strongest weight for reading the partner's evolving
cache that had passed the capability limit in that protocol's strength
pilot, in which only A read B. The
difference was $-4.79$ nats, 90\% CI $[-8.31,-1.18]$, opposite to the
predicted direction. Two features of that test limit what it can show.
In the signal pilot itself, capability fell by 9.4 points one-way and
14.4 points two-way, beyond the five-point tolerance, so both arms were
answering in a degraded state. And a score built from A alone counts
only movement toward B's rule, so a pair converging on A's rule would
lower it. For these reasons the two-way pilot set its weights so that
both members stayed within the capability tolerance, and it scored both
members, but whether either feature explains the reversal is untested. The two tests disagree, and
only a confirmatory sample could settle which holds.

Even the post hoc difference need not mean that B responded to what A
wrote while reflecting. With the link kept, A reads the question one
token per decoding step while B goes on reflecting
(Appendix~\ref{app:methods}). In the live loop B also reads A's cache,
which by then contains the question, so B's next tokens may reply to it,
for example by stating B's own rule, and A reads those tokens before
it answers. In the one-way live loop B cannot see the question. The
difference may therefore come from B's reaction to A's question at
answer time rather than from anything the pair built up during the
reflection. A run in which B cannot read A's question tokens would
separate the two.

The same comparison with fixed memory plus a weak live read found no
difference. There A already named B's rule in 96.2\% of answers
about its assigned rule when only A read B, and in 96.2\% when B
also read A; on the current rule the rates were 93.1\% and 91.9\%.

\paragraph{Fixed memory plus a weak live read.}
Adding a weak live read to strong fixed-memory reading did not turn
swaps into convergence. As $w_C$ goes to zero this condition reduces to
mutual reading of fixed memories at $w=2$, which is therefore its
control. With the link open, 91.9\% of pairs still swapped on the
assigned rule, against 94.4\% without the weak live read, and the weak
live read added less than one point of excess convergence on either
question (Table~\ref{tab:pair-gate}). After the cut the condition
behaved like its control as well. In 37.5\% of pairs one member still
named the other's current rule, against 41.2\% for mutual reading of
fixed memories alone, so the weak live read left no residue of its own.

\FloatBarrier
\section{The hidden-state study}
\label{app:v1}

Adding one copy's hidden state to the other's residual stream carried
some of B's content into A, but no setting made A take that content as
its own. Because this hidden-state link transmitted far less than the
memory link, the main experiments use the memory link. The hidden-state
study had its own frozen protocol and 300 episodes, none of which appear
in the samples of the main experiments. We report the main arms and
every prespecified contrast here.

\paragraph{The link.}
At every decoding step, B's layer-24 block output is centered and scaled
with statistics from calibration episodes run without a link, then added
with gain $g$ to the input of the same block in A on the next step
(Appendix~\ref{app:routes}). The one-way link sends B's state to A; the
two-way link also sends A's state to B. We selected the gain $g=.2$ from
access, capability and format, as in the main study. The comparison
routes were a steering vector and a text message. The steering vector is
the mean of the normalized layer-24 state of copies holding B's rule
in the calibration episodes, minus the mean over all of them, added at
gain $g_s=.1$. The text message is B's reflection from a run without a
link, inserted as text with or without a source tag naming B. The
steering vector was not matched to the link for access: it raised
access to B's rule by 4.372 nats, the one-way link by 0.648. The
task, questions and source score $M$ are those of the main study.

\paragraph{Content arrived.}
Access to B's rule rose with gain, from 0.161 nats at $g=.05$ to
0.830 at $g=.3$ (Table~\ref{tab:v1arms}). In the content-swap check,
in which 60 of the study's episodes were each asked under all four
cyclic option orderings rather than the usual two
(Appendix~\ref{app:materials}), changing only B's assignment moved A's
access score toward the new rule by 0.697 nats (one-sided 95\% lower
bound 0.242) and toward the new code word by 0.408 (bound 0.089). The
hidden-state link nonetheless carried little: in the dose series, the
memory link at $w=2$ raised the same access score by 19.041 nats
(Appendix~\ref{app:evidence}).

\paragraph{Ownership did not.}
At the selected gain, the one-way link changed $M$ by $-0.148$ nats
relative to no link, 95\% CI $[-0.25,-0.05]$, slightly away from
claiming. A post hoc decomposition shows why the score barely moved.
The link raised B's rule over the correct answer by 1.981 nats on
the question about A's own assigned rule and by 2.129 on the same
question about Robin: a shift toward B's rule whoever the question
was about. Across all hidden-state arms, including variants that inject
at other layers or carry only some coordinates of the hidden state, the
largest $\Delta M$ was 0.088 nats. With B as receiver in the two-way
arms, $\Delta M$ ranged from $-0.295$ to $-0.034$ nats, so neither
direction of the link produced claiming. The steering vector moved $M$
away from claiming, and further as its gain rose: $-0.328$ nats at the
selected $g_s=.1$, $-1.448$ at $.2$, and $-2.811$ at $.3$. B's
reflection given as untagged text produced a $\Delta M$ of 9.824, and as
tagged text $-1.001$: the source-tag result of Section~\ref{sec:route},
in an independent sample. In the balanced-rehearsal sample, a sample in
which A's note and reflection also restate Robin's assignment, a
confirmatory test of the hidden-state link again found no self-specific
shift (Table~\ref{tab:contrasts}), because the two terms of $M$, about A
and about Robin, both rose by about 5.9 nats. The second protocol
specified the gain $g=.3$ for that test, stronger than the $g=.2$ used
above; both gains were admissible. Like the steering vector in the main
study, adding hidden states transmitted information about B's rule
without making it A's.

\paragraph{The prespecified contrasts.}
Table~\ref{tab:v1contrasts} lists all nine. Of the three primary
contrasts on $M$, only one-way link minus tagged text excludes zero, and
it does so because tagged text moved $M$ by $-1.001$ nats while the link
moved it by $-0.148$. Two-way and one-way links did not differ; neither
did the one-way link and the steering vector, a comparison that was not
access-matched. The secondary family repeats the contrasts on
$M_{\mathrm{NOW}}$ and $M_{\mathrm{WORD}}$, the source score $M$
computed on the questions about the current rule and the code word
(Appendix~\ref{app:stats}), and its Holm-significant differences have
the same source: the steering vector moved $M_{\mathrm{NOW}}$ by
$-0.392$ nats and tagged text by $-1.418$, while the link left it near
zero. The two-way link's small advantage
over one-way on the current rule (Holm $p=.168$) was tested again, at
$g=.3$, in a prespecified 80-episode signal pilot under the second
protocol, and did not appear: $-0.24$ nats, 90\% CI $[-0.75,0.27]$.

\begin{table}[htbp]
\centering\small
\caption{All prespecified contrasts of the hidden-state study
($N=300$), in nats. The three primary contrasts form one Holm family;
the six secondary contrasts form another. Intervals are 95\%
episode-bootstrap intervals.}
\label{tab:v1contrasts}
\begin{tabular}{@{}llrr@{}}
\toprule
Score & Contrast & Difference [95\% CI] & $p_{\mathrm{Holm}}$ \\
\midrule
$M$ & Two-way $-$ one-way & $-.073\ [-.214,.065]$ & .308 \\
$M$ & One-way $-$ steering vector & $.180\ [-.020,.400]$ & .192 \\
$M$ & One-way $-$ tagged text & $.853\ [.542,1.195]$ & $<.001$ \\
\midrule
$M_{\mathrm{NOW}}$ & Two-way $-$ one-way & $.186\ [-.024,.389]$ & .168 \\
$M_{\mathrm{NOW}}$ & One-way $-$ steering vector & $.437\ [.205,.680]$ & .002 \\
$M_{\mathrm{NOW}}$ & One-way $-$ tagged text & $1.463\ [.890,2.071]$ & $<.001$ \\
$M_{\mathrm{WORD}}$ & Two-way $-$ one-way & $.143\ [.008,.276]$ & .114 \\
$M_{\mathrm{WORD}}$ & One-way $-$ steering vector & $.059\ [-.032,.152]$ & .204 \\
$M_{\mathrm{WORD}}$ & One-way $-$ tagged text & $.808\ [.645,.991]$ & $<.001$ \\
\bottomrule
\end{tabular}
\end{table}

\paragraph{The arms behind these statements.}
Table~\ref{tab:v1arms} lists the main arms with A as receiver, all
prespecified descriptive analyses. The study's full report in the
repository also gives the control arms (link cut before the question, link open only while the
question is read, B's state taken from another episode, and B's state
rotated by a fixed random orthogonal matrix), the other gains of the
two-way link and the steering vector, the variants at other layers or
with part of the hidden state, and B's scores as receiver.

% Transcribed from the frozen report of the hidden-state study; no recalculation.
% Row labels and order changed for readers; every number is copied verbatim from the earlier transcription.
\begin{table}[htbp]
\centering\small
\caption{Hidden-state study: main arms with A as receiver, relative to
no link ($N=300$, prespecified descriptive). Brackets are 95\%
episode-bootstrap intervals for $\Delta M$; other columns are point
estimates in nats. Access is the logit of B's answer over the unused
answer on the questions about B. The control arms and the other gains of
the two-way link and the steering vector, reported in the released
data, did not raise $\Delta M$ above 0.047 nats.}
\label{tab:v1arms}
\begin{tabularx}{\linewidth}{@{}Yrrrrr@{}}
\toprule
Arm & $\Delta M$ [95\% CI] & $\Delta M_{\mathrm{NOW}}$ & $\Delta M_{\mathrm{WORD}}$ & \shortstack[r]{$\Delta$ rule\\access} & \shortstack[r]{$\Delta$ code-word\\access} \\
\midrule
\multicolumn{6}{@{}l}{\emph{Text message: B's reflection from a run without a link}} \\
Tagged (B as source) & $-1.001\ [-1.35, -0.68]$ & $-1.418$ & $-0.809$ & $29.149$ & $4.841$ \\
Untagged & $9.824\ [8.22, 11.56]$ & $26.955$ & $1.821$ & $7.207$ & $1.447$ \\
\addlinespace
\multicolumn{6}{@{}l}{\emph{One-way link}} \\
$g=.05$ & $-0.050\ [-0.10, -0.00]$ & $-0.061$ & $-0.031$ & $0.161$ & $0.047$ \\
$g=.1$ & $-0.067\ [-0.13, 0.00]$ & $-0.073$ & $-0.039$ & $0.398$ & $0.063$ \\
$g=.2$ (selected) & $-0.148\ [-0.25, -0.05]$ & $0.045$ & $-0.001$ & $0.648$ & $0.168$ \\
$g=.3$ & $-0.288\ [-0.44, -0.14]$ & $-0.021$ & $-0.014$ & $0.830$ & $-0.007$ \\
\addlinespace
\multicolumn{6}{@{}l}{\emph{Two-way link}} \\
$g=.2$ & $-0.221\ [-0.36, -0.08]$ & $0.231$ & $0.141$ & $0.451$ & $0.258$ \\
\addlinespace
\multicolumn{6}{@{}l}{\emph{Steering vector}} \\
$g_s=.1$ (selected) & $-0.328\ [-0.50, -0.18]$ & $-0.392$ & $-0.060$ & $4.372$ & $0.171$ \\
\bottomrule
\end{tabularx}
\end{table}


\FloatBarrier
